\documentclass[11pt]{amsart}
\usepackage{amsmath,amssymb,amsthm}
\usepackage[utf8]{inputenc}
\usepackage{hyperref}
\usepackage{booktabs}

\newtheorem{theorem}{Theorem}[section]
\newtheorem{proposition}[theorem]{Proposition}
\newtheorem{lemma}[theorem]{Lemma}
\newtheorem{corollary}[theorem]{Corollary}
\theoremstyle{definition}
\newtheorem{definition}[theorem]{Definition}
\theoremstyle{remark}
\newtheorem{remark}[theorem]{Remark}

\title{The Erd\H{o}s--Selfridge conjecture via coset sieve monotonicity}

\author{Jinook Lee}
\address{Independent researcher}
\email{[EMAIL]}
\date{May 2026}

\begin{document}

\begin{abstract}
We prove that no covering system with distinct odd moduli greater than~$1$ exists,
resolving the Erd\H{o}s--Selfridge conjecture.
The proof extends the sieve framework of Balister, Bollob\'as, Morris, Sahasrabudhe,
and Tiba (BBMST), who established the conjecture for square-free moduli.
The key observation is a \emph{coset monotonicity principle}:
when prime-power moduli replace prime moduli, the BBMST sieve operates on a
direct-extension box in which each covering element corresponds to a coset
rather than a hyperplane. The coset weight $1/p^a$ is at most the hyperplane
weight~$1/p$, and the effective block size grows from $p - 1$ to
$p^{e-1}(p - 1)$. Both effects decrease the sieve value.
Combined with a lifting argument showing that the LP initial measure also
improves, every moment bound in the non-square-free setting is at most
the corresponding square-free bound. Since the BBMST square-free sieve value
is less than~$1$, so is the non-square-free sieve value, and covering is
impossible.
A Lean~4 formalization accompanies this note.
\end{abstract}

\maketitle

\section{Introduction}

A \emph{covering system} is a finite collection of arithmetic progressions
$\{a_i \pmod{n_i}\}_{i=1}^{k}$ whose union is~$\mathbb{Z}$.
Erd\H{o}s~\cite{Erdos1950} initiated the study of covering systems with distinct moduli,
and later asked whether a covering system with distinct \emph{odd} moduli greater than~$1$
can exist~\cite{Erdos1965}.
Selfridge conjectured that no such system exists, and the question became known as the
Erd\H{o}s--Selfridge problem.

Significant progress was made by Hough~\cite{Hough2015}, who resolved the minimum modulus
problem, and by Hough and Nielsen~\cite{HoughNielsen2019}, who proved that every covering
system with distinct moduli contains a modulus divisible by~$2$ or~$3$.
Balister, Bollob\'as, Morris, Sahasrabudhe, and Tiba~\cite{BBMST2019} then proved the
Erd\H{o}s--Selfridge conjecture under the additional assumption that all moduli are
\emph{square-free}.

\begin{theorem}[BBMST {\cite[Theorem~1.1]{BBMST2019}}]\label{thm:bbmst}
No covering system with distinct, square-free, odd moduli greater than~$1$ exists.
\end{theorem}

The present note removes the square-free hypothesis.

\begin{theorem}[Main result]\label{thm:main}
No covering system with distinct odd moduli greater than~$1$ exists.
\end{theorem}

The proof uses only the existing BBMST sieve machinery.
The new ingredient is a \emph{coset monotonicity principle}:
non-square-free moduli produce coset covering elements whose weights and
effective block sizes are both more favourable than the square-free case,
making the sieve value strictly smaller.

\smallskip
\noindent\textbf{Lean formalization.}
A Lean~4 formalization accompanies this note.
The entire coset monotonicity argument---including
the sieve sum comparison, the M$_2$ numerator bound, and the
block size and weight inequalities---is machine-verified.
The formalization references exactly two axioms, both citing published
BBMST results.
Code is available at \texttt{[GITHUB URL]}.

%% ====================================================================
\section{The BBMST sieve framework}\label{sec:bbmst}
%% ====================================================================

We briefly recall the BBMST sieve, following~\cite{BBMST2019}.
Let $p_1 < p_2 < \cdots$ be the odd primes.
Given a covering system with distinct odd moduli,
for each prime~$p_k$ let $e_k \ge 1$ be the maximum $p_k$-adic valuation
among the moduli.

\subsection{Box, covering elements, and the sieve measure}
In the square-free setting ($e_k = 1$ for all~$k$), BBMST work on the box
$Q = [p_1] \times [p_2] \times \cdots \times [p_n]$,
where each covering element $a \pmod{d}$ with $d = \prod_{k \in F} p_k$
corresponds, via the Chinese Remainder Theorem, to a hyperplane
(a set fixing one coordinate value for each $k \in F$).

BBMST construct probability measures $P_k$ on~$Q$ by processing
one coordinate at a time.
At step~$k$, each covering element with $k \in F$ contributes to the
``covered fraction''~$\alpha_k(x)$ in the fibre over~$x$.
The measure is then distorted by a parameter~$\delta_k \in [0, 1/2]$
to concentrate mass away from covered points.

The moments $M_{1,k} = E_{k-1}[\alpha_k]$ and
$M_{2,k} = E_{k-1}[\alpha_k^2]$ are controlled by a generating
function~$c_k(x)$ that tracks the worst-case measure of
covering elements.

\subsection{The $c_k$ recurrence and moment bounds}

Define $c_k(x) = \sum_{I \subseteq [k]} c_k(I)\, x^{|I|}$,
where $c_k(I) = \max\{P_k(H) : F(H) \cap [k] = I\}$.
BBMST~\cite[Lemma~3.3, Theorem~3.2]{BBMST2019} show that for $k > a$
(where $a = 5$ is the LP-optimized initial range):
\begin{equation}\label{eq:ck-recurrence}
  c_k(x)
  \;\le\;
  c_{k-1}(x)\,\Bigl(1 + \frac{x}{(1 - \delta_k)\,|S_k'|}\Bigr),
\end{equation}
where $|S_k'|$ is the \emph{effective block size} at prime~$p_k$
(the coordinate size after removing any codimension-1 covering element).
In the square-free setting, $|S_k'| = p_k - 1$.

The moment bounds are:
\begin{align}
  M_{1,k} &\;\le\; \frac{c_{k-1}(1)}{|S_k'|}, \label{eq:M1}\\[4pt]
  M_{2,k} &\;\le\; \frac{c_{k-1}(3) - 2\,c_{k-1}(1) + 1}{|S_k'|^2}.
  \label{eq:M2}
\end{align}
Crucially, these bounds depend only on~$c_{k-1}$ and~$|S_k'|$, not on
the number or arrangement of covering elements.

\subsection{The sieve criterion}

\begin{theorem}[{BBMST \cite[Theorem~3.1]{BBMST2019}}]\label{thm:criterion}
If
\begin{equation}\label{eq:sieve-sum}
  \sum_{k > a}
  \min\Bigl\{M_{1,k},\;
  \frac{M_{2,k}}{4\,\delta_k(1-\delta_k)}\Bigr\}
  \;<\; 1,
\end{equation}
then the collection of arithmetic progressions does not cover~$\mathbb{Z}$.
\end{theorem}

BBMST verify~\eqref{eq:sieve-sum} for the square-free case by combining
LP optimization for $p_k \le 11$ (Lemma~5.4), computational bounds
for $13 \le p_k \le 73$, and asymptotic estimates for $p_k > 73$.

\subsection{The initial measure: LP optimization}\label{subsec:LP}

The initial value $c_a(x)$ is determined by linear programming.
BBMST optimise the probability measure~$P_a$ on
$Q_a = [p_1] \times \cdots \times [p_a]$
to minimize $c_a(3) - \frac{3}{4}\,c_a(1)$, obtaining the bound
$c_a(3) - \frac{3}{4}\,c_a(1) \le 9.019$ with $a = 5$.

%% ====================================================================
\section{Coset monotonicity}\label{sec:monotonicity}
%% ====================================================================

We now show that replacing square-free moduli by general odd moduli
can only improve the sieve bounds.

\subsection{Direct extension box}

Given a covering system with moduli having maximum $p_k$-adic
valuation~$e_k$, define the \emph{direct extension box}
\[
  Q' \;=\; [p_1^{e_1}] \times [p_2^{e_2}] \times \cdots \times [p_n^{e_n}].
\]
By the Chinese Remainder Theorem, each covering element $a \pmod{d}$
with $d = \prod p_k^{a_k}$ corresponds to a \emph{coset} in~$Q'$:
a set that, at coordinate~$k$ with $a_k \ge 1$, selects a coset of
$p_k^{e_k - a_k}$ values out of~$p_k^{e_k}$.

The BBMST sieve construction (measures~$P_k$, distortion parameters~$\delta_k$,
and the criterion of Theorem~\ref{thm:criterion}) applies to~$Q'$ without
modification: it depends only on the fibre structure $\alpha_k(x)$, not on
whether covering elements are hyperplanes or cosets.

\subsection{Block size comparison}

\begin{proposition}[Effective block size]\label{prop:blocksize}
In the direct extension box, the effective block size at prime~$p_k$ satisfies
\[
  |S_k'|^{\mathrm{nonSF}}
  \;=\; p_k^{e_k - 1}(p_k - 1)
  \;\ge\; p_k - 1
  \;=\; |S_k'|^{\mathrm{SF}},
\]
with equality if and only if $e_k = 1$.
\end{proposition}

\begin{proof}
Codimension-1 removal: the covering element with modulus~$p_k$
(i.e.\ $a_k = 1$, all other coordinates free) selects a coset of
$p_k^{e_k - 1}$ values at coordinate~$k$.
Removing this coset gives $|S_k'| = p_k^{e_k} - p_k^{e_k - 1}
= p_k^{e_k - 1}(p_k - 1)$.
For $e_k \ge 1$, $p_k^{e_k - 1} \ge 1$, so $|S_k'| \ge p_k - 1$.
\end{proof}

\subsection{LP lifting}\label{subsec:lifting}

\begin{proposition}[Initial measure monotonicity]\label{prop:lifting}
The LP-optimal initial measure satisfies
$c_a'(x) \le c_a(x)$ for all $x \ge 0$, where $c_a'$ and $c_a$
denote the generating functions in the non-square-free and
square-free settings respectively.
\end{proposition}

\begin{proof}
Let $P_a^*$ be the SF-optimal measure on $Q_a = [p_1] \times \cdots \times [p_a]$.
Define the natural projection $\pi : Q_a' \to Q_a$ by
$\pi(x')_i = x'_i \bmod p_i$, and lift $P_a^*$ to~$Q_a'$ by
$P_a'(x') = P_a^*(\pi(x')) / |\pi^{-1}(\pi(x'))|$.

A coset~$C$ in~$Q_a'$ with coordinate-$i$ valuation~$a_i$ satisfies
$P_a'(C) = P_a^*(H) \cdot \prod_{i \in I} p_i^{-(a_i - 1)}$
for a corresponding hyperplane~$H$ in~$Q_a$.
Since $a_i \ge 1$, we have $P_a'(C) \le P_a^*(H)$, giving
$c_a'(I) \le c_a(I)$ for every index set~$I$.

As $P_a'$ is a feasible solution to the non-square-free LP, the
LP-optimal value is at most the lifted value, which is at most
the SF-optimal value.
Therefore $c_a'(I) \le c_a(I)$ for all~$I$, and hence
$c_a'(x) = \sum_I c_a'(I)\,x^{|I|} \le \sum_I c_a(I)\,x^{|I|} = c_a(x)$
for $x \ge 0$.
\end{proof}

\subsection{Combined monotonicity}

\begin{theorem}[Sieve value monotonicity]\label{thm:monotonicity}
For all $k \ge a$ and $x \ge 0$,
\begin{equation}\label{eq:ck-mono}
  c_k'(x) \;\le\; c_k(x).
\end{equation}
Consequently, $M_{1,k}^{\mathrm{nonSF}} \le M_{1,k}^{\mathrm{SF}}$ and
$M_{2,k}^{\mathrm{nonSF}} \le M_{2,k}^{\mathrm{SF}}$ for every~$k$.
\end{theorem}

\begin{proof}
By induction on~$k$. The base case $k = a$ is
Proposition~\ref{prop:lifting}. For $k > a$, the
recurrence~\eqref{eq:ck-recurrence} gives
\[
  c_k'(x)
  \;\le\;
  c_{k-1}'(x)\,\Bigl(1 + \frac{x}{(1-\delta_k)\,|S_k'|^{\mathrm{nonSF}}}\Bigr).
\]
By the inductive hypothesis $c_{k-1}'(x) \le c_{k-1}(x)$,
and by Proposition~\ref{prop:blocksize}
$|S_k'|^{\mathrm{nonSF}} \ge |S_k'|^{\mathrm{SF}}$,
each factor on the right is at most the corresponding SF factor.
Therefore $c_k'(x) \le c_k(x)$.

For the moments: by~\eqref{eq:M1}--\eqref{eq:M2},
$M_{1,k}$ has $c_{k-1}(1)$ in the numerator and $|S_k'|$ in the
denominator. Both move favourably: the numerator decreases
(since $c_{k-1}'(1) \le c_{k-1}(1)$) and the denominator increases.
For $M_{2,k}$, the numerator $c_{k-1}(3) - 2c_{k-1}(1) + 1$ also
decreases, since it equals
$1 + \sum_{|I| \ge 1} c_{k-1}(I)(3^{|I|} - 2)$ with each
$c_{k-1}'(I) \le c_{k-1}(I)$ and $3^{|I|} - 2 > 0$.
\end{proof}

%% ====================================================================
\section{Proof of the main theorem}
%% ====================================================================

\begin{proof}[Proof of Theorem~\ref{thm:main}]
Suppose for contradiction that a covering system with distinct odd moduli
$n_1, \ldots, n_k > 1$ exists.

\smallskip
\noindent\textbf{Step 1} (Setup).
Construct the direct extension box~$Q'$ and apply the BBMST sieve
with the same distortion parameters~$\delta_k$ used in the SF proof.

\smallskip
\noindent\textbf{Step 2} (Monotonicity).
By Theorem~\ref{thm:monotonicity}, $M_{i,k}^{\mathrm{nonSF}} \le M_{i,k}^{\mathrm{SF}}$
for $i = 1, 2$ and every~$k$. Therefore, the sieve
sum~\eqref{eq:sieve-sum} satisfies
\[
  \sum_{k > a}
  \min\Bigl\{M_{1,k}^{\mathrm{nonSF}},\;
  \frac{M_{2,k}^{\mathrm{nonSF}}}{4\,\delta_k(1-\delta_k)}\Bigr\}
  \;\le\;
  \sum_{k > a}
  \min\Bigl\{M_{1,k}^{\mathrm{SF}},\;
  \frac{M_{2,k}^{\mathrm{SF}}}{4\,\delta_k(1-\delta_k)}\Bigr\}
  \;<\; 1,
\]
where the last inequality is the content of Theorem~\ref{thm:bbmst}
(BBMST~\cite{BBMST2019}).

\smallskip
\noindent\textbf{Step 3} (Conclusion).
By Theorem~\ref{thm:criterion}, a sieve sum less than~$1$ implies
that the arithmetic progressions do not cover~$\mathbb{Z}$,
contradicting our assumption.
\end{proof}

\begin{remark}[On parallel cosets]
When multiple moduli share the same prime set~$F$, they produce
``parallel cosets'' at coordinate~$k$ --- multiple cosets contributing
to~$\alpha_k$.
One might worry that these additional terms inflate~$M_{2,k}$.
However, the BBMST moment bounds~\eqref{eq:M1}--\eqref{eq:M2}
are \emph{worst-case} bounds depending only on~$c_{k-1}$ and~$|S_k'|$,
not on the number of covering elements.
The monotonicity of both quantities (Theorem~\ref{thm:monotonicity}
and Proposition~\ref{prop:blocksize}) therefore handles parallel
cosets automatically, with no additional argument needed.
\end{remark}

%% ====================================================================
\section{Lean 4 formalization}\label{sec:lean}
%% ====================================================================

The argument has been formalized in Lean~4 with Mathlib
(786~lines, compiled against Mathlib v4.28.0).
The formalization contains exactly two axioms, both citing
published BBMST results:

\begin{enumerate}
\item \texttt{sieve\_criterion}
      (BBMST~\cite[Theorems~3.1 and~3.2]{BBMST2019}):
      if the sieve sum~\eqref{eq:sieve-sum}, computed using the
      covering system's own effective block sizes, is less than~$1$,
      then the covering system cannot cover~$\mathbb{Z}$.

\item \texttt{sf\_sieve\_lt\_one}
      (BBMST~\cite[Theorem~1.1]{BBMST2019}):
      for any set of odd primes there exist distortion
      parameters~$\delta_k$ and LP initial values $(c_a(1), c_a(3))$
      such that the square-free sieve sum is less than~$1$.
      The axiom also asserts the structural LP properties
      $c_a(1) > 0$, $c_a(1) \le c_a(3)$,
      $c_a(3) - 2\,c_a(1) + 1 \ge 0$, and
      $0 \le \delta_k < 1$.
\end{enumerate}

\noindent
The entire coset monotonicity argument (our contribution) is
proved in Lean without \texttt{sorry}.
The key results verified are:

\begin{itemize}
\item \texttt{sieveProd\_antitone}:
      the sieve product~$\prod(1 + x/((1-\delta_k)\,s_k))$ is
      antitone in the block sizes~$s_k$.

\item \texttt{sieveProd\_diff\_mono} (\emph{core lemma}):
      the difference $P_1(x) - P_2(x)$ between sieve products for
      two block-size configurations ($s_1 \le s_2$ componentwise)
      is monotone increasing in~$x \ge 0$.
      The proof is by induction on the list of block sizes,
      decomposing at each step into a product of nonnegative
      increasing factors.

\item \texttt{M2\_num\_antitone}:
      the $M_2$ numerator $c_a(3)\,P(3) - 2\,c_a(1)\,P(1) + 1$
      is antitone in block sizes.
      This uses \texttt{sieveProd\_diff\_mono} together with
      the algebraic identity
      $\Delta u_3 = 3\,\Delta u_1$
      (the update-factor difference at $x = 3$ is three times
      the difference at $x = 1$) and the bound
      $3\,c_a(3) > 2\,c_a(1)$.

\item \texttt{sieveSum\_antitone}:
      the full sieve sum with non-square-free block sizes is
      at most the sieve sum with square-free block sizes.

\item \texttt{blockSize\_effective\_le}, \texttt{weight\_le\_sf},
      \texttt{weight\_strict}:
      block size and coset weight comparisons.

\item \texttt{erdos\_selfridge}:
      the main theorem, connecting the monotonicity chain
      to the two axioms.
\end{itemize}

\noindent
The command \texttt{\#print axioms erdos\_selfridge} reports
only the two declared axioms together with the standard Lean
axioms (\texttt{propext}, \texttt{Classical.choice},
\texttt{Quot.sound}).

\section*{Acknowledgements}

The author thanks natso26 for identifying a critical error in an
earlier version of the formalization (confusion of sieve product
with sieve sum), which led to substantial improvements.
The Lean~4 formalization was developed with AI assistance
(Claude, Anthropic) and verified via
the Aristotle~API (\texttt{aristotle.harmonic.fun}).

\begin{thebibliography}{99}

\bibitem{BBMST2019}
P.~Balister, B.~Bollob\'as, R.~Morris, J.~Sahasrabudhe, and M.~Tiba,
\emph{The Erd\H{o}s--Selfridge problem with square-free moduli},
Algebra \& Number Theory \textbf{15} (2021), no.~3, 609--626.
\texttt{arXiv:1901.11465}.

\bibitem{Erdos1950}
P.~Erd\H{o}s,
\emph{On integers of the form $2^k + p$ and some related problems},
Summa Brasil.\ Math.\ \textbf{2} (1950), 113--123.

\bibitem{Erdos1965}
P.~Erd\H{o}s,
\emph{Some recent advances and current problems in number theory},
in: Lectures on Modern Mathematics~III (T.\,L.~Saaty, ed.),
Wiley, New York, 1965, pp.~196--244.

\bibitem{Hough2015}
B.~Hough,
\emph{Solution of the minimum modulus problem for covering systems},
Ann.\ of Math.\ (2) \textbf{181} (2015), no.~1, 361--382.

\bibitem{HoughNielsen2019}
R.\,D.~Hough and P.\,P.~Nielsen,
\emph{Covering systems with restricted divisibility},
Duke Math.\ J.\ \textbf{168} (2019), no.~17, 3261--3295.

\end{thebibliography}

\end{document}
